\documentclass[preprint,12pt]{elsarticle}
\usepackage[T1]{fontenc}
\usepackage{lmodern}
\usepackage{amsmath,amssymb,booktabs}
\usepackage{array}
\usepackage{algorithm}
\usepackage{algpseudocode}
\usepackage{capt-of}
\usepackage{listings}
\usepackage{xcolor}
\usepackage{natbib}
\usepackage{xurl}
\usepackage{hyperref}
\def\algorithmautorefname{Algorithm}
\journal{Journal of Systems and Software}
\myfooter[C]{Author-review draft for the Journal of Systems and Software}
\definecolor{ealKeyword}{RGB}{31,73,112}
\definecolor{ealString}{RGB}{112,53,29}
\definecolor{ealComment}{RGB}{73,91,80}
\lstdefinelanguage{EALTwo}{
  morekeywords={language,environment,require,tool,version,mode,deterministic,
    nondeterministic,evidence,kind,max_age,input,assumption,statement,validate,
    valid_from,valid_until,reasoning,method,rationale,backing,claim,proposition,
    subject,quantity,unit,scope,query,result,argument,conclusion,premises,
    assumptions,binding,objection,target,pattern,apply,true,false,null},
  sensitive=true,morecomment=[l]{//},morecomment=[s]{/*}{*/},morestring=[b]"
}
\lstset{language=EALTwo,basicstyle=\ttfamily\footnotesize,
  keywordstyle=\color{ealKeyword}\bfseries,stringstyle=\color{ealString},
  commentstyle=\color{ealComment}\itshape,
  columns=fullflexible,keepspaces=true,breaklines=true,
  breakatwhitespace=true,showstringspaces=false,
  aboveskip=0.7\baselineskip,belowskip=0.7\baselineskip,captionpos=b}

\begin{document}
\begin{frontmatter}
\title{EAL/2: Executable engineering arguments for critical review of AI-assisted decisions}
\author[author]{Emmett Miller}
\address[author]{Affiliation and corresponding-author contact pending author confirmation}
\begin{abstract}
EAL/2 makes claims, evidence, methods and objections explicit so engineers can challenge AI advice. We describe its evaluator and claim-scoped host and examine retained synthetic experiments. An earlier raw-only diagnostic scored EAL 4/10 versus meaning-equivalent JSON 5/10 on exposed tasks. An amended 96-call negative-revision study accepted no generated source under strict validation. Among 960 mechanism assignments, 460 matched the supplied-information reference and 130 admitted false support. An amended 800-slot deployment used 804 requests: direct interpretation gave 50/144 exact statuses, while host and native routes retained 240/240 checked statuses by construction. All 32 final authored sources failed validation; 320 recipient packets were unavailable. A separate bias-attribution Stage A run reached 432 terminal call outcomes, of which 160 passed response validation. Five of 288 model--case--route dispositions met the full author-labelled reference. Twenty strong false attributions remained in parsed analyst answers rejected by subsequent validation. Descriptive family-bootstrap topology intervals all included zero; the one-call direct and two-call independent routes had unequal allocation. The execution rule changed after the first invalid pilot response, so the run is developmental. These author-created cases reveal response, source-authoring and communication failures while leaving notation benefit, general model ranking and human bias mitigation unestablished. Fresh-family representation and topology comparisons and a randomised human-decision trial remain prospective. Independent adjudication and complete effort accounting are needed to evaluate practical benefit.
\end{abstract}
\begin{keyword}
engineering argumentation \sep domain-specific language \sep defeasible reasoning \sep language models \sep cognitive bias \sep human--AI decision making
\end{keyword}
\end{frontmatter}

\section{Introduction}\label{sec:introduction}
The motivating concern behind EAL/2 is a human review problem: as an engineering team delegates analysis or drafting to an AI agent, a fluent answer may be accepted without checking whether its observations answer the engineering question. In a controlled non-AI task, easy-to-read statements received higher truth ratings \citep{reber1999}; another experiment reversed the direction after participants learned that fluency predicted falsehood \citep{unkelbach2007}. In human--AI studies, explanations increased acceptance of both correct and incorrect advice \citep{bansal2021}, and longer LLM explanations raised confidence without improving discrimination between correct and incorrect answers \citep{steyvers2025}. These results motivate a context-sensitive test in engineering review; they do not show that a particular EAL/2 user's decision was biased. An engineering decision can turn on a detail that a plausible summary omits. A vibration record may show a value below a limit, yet its calculation can use a different origin from the claim's formal question. A challenge to an argument can depend on the very claim it challenges, leaving the dispute unresolved until independent grounds are supplied. A status label without its dependencies is difficult to review or revise.

EAL/2 recasts the reviewable part of an engineering problem as arguments: an explicit question, grounds, inference method, assumptions, objections and a claim whose scope can be checked. The original purpose was to support critical analysis by a person reviewing AI output. Its evaluator checks supplied evidence and method outputs, then computes the status of the resulting support--attack structure. A text model may author or interpret that structure, but its fluent account does not perform the evaluator's checks. The engineering question is therefore twofold: what can the language and evaluator check, and which parts of their result survive transfer to a model recipient? The Stage A study asks whether agents using that structure can match stipulated process records to bounded bias dispositions while retaining technical and organisational rivals.

This article contributes a precise account of EAL/2's implemented profile, a worked case that distinguishes numerical success from a matching formal query, and an analysis of retained model experiments. The relay comparison tests hand-offs from an already correct tool-enabled producer. The notation diagnostic introduces wrong proposed answers and a paired EAL--JSON comparison. Later recipient pilots supplied a checked status to one model on selected synthetic states. The newer host and reviewed-family implementation changes the delivery boundary and motivates the separate cross-model comparison reported below. The later 96-call developmental schedule includes linked model drafting and revision under a strict checker, which accepted none of its generated submissions. The original version 0.1.0 protocols remain unrun: their human adjudication and broader task sampling were not supplied by the amended synthetic schedules. Neither an advantage caused by EAL syntax nor reliable model preservation of qualifications has been shown. The completed developmental Stage A study in Section~\ref{sec:bias-stage-a} uses a separate process-record reference. Its 432 terminal calls support a reproducible failure audit of synthetic bias dispositions; they contain no measured human decision benefit.

\section{Language and executable meaning}\label{sec:language}
\subsection{Design boundary}
EAL/2 is the repository's sole supported source language. Its grammar requires the header \texttt{language "EAL/2";}; earlier EAL measurements retain their historical versions and cannot be relabelled as EAL/2 results. The ANTLR grammar recognises source, which is lowered and checked against a typed internal representation before evaluation. A \texttt{reasoning} declaration selects one versioned method, for example \texttt{engineering/rms/1}. The tool's \texttt{mode} instead describes collection variability. Built-in and installed methods use the same contract and result-binding checks. Reusable patterns have typed parameters, closed lexical scope and a finite expansion into ordinary arguments; expansion retains the identity of shared observations \citep{earlDesign}.

The language makes an authored inference reviewable, rather than turning prose into a proof. A claim has a statement and an environment. Evidence declares its acquisition tool, kind, scope, age limit and value predicates. Assumptions add validation evidence and optional time bounds. An argument names a conclusion, a reasoning step and the evidence, assumptions or premise claims on which that step depends. An objection names the exact claim, argument, method, assumption or other objection it challenges. Toulmin's separation of grounds, inference and rebuttal informs this vocabulary \citep{toulmin2003}. SACM defines a metamodel and graphical notation for structured assurance cases \citep{sacm2021}. EAL/2 specifies executable method contracts and status evaluation for its own bounded engineering-argument profile; neither specification's existence establishes that one is easier to use. EAL/2 does not implement the complete argument-construction, contrariness or preference machinery of ASPIC+ \citep{modgil2014}.

\subsection{A query that a numerical threshold cannot settle}
Listing~\ref{lst:rms} is a complete, existing EAL/2 programme from the repository. It asks whether recorded root-mean-square (RMS) pressure about zero is at most \(5\ \mathrm{kPa}\) for a named sample episode. The external method is identified by a versioned contract; the claim's proposition fixes the subject, quantity, unit, scope, time and query. The \texttt{binding} clause requires the observation used by the method to answer that formal query.
\begin{lstlisting}[caption={Complete EAL/2 RMS example; reproduced from examples/rms.eal at the repository revision stated in the data availability section.},label={lst:rms}]
language "EAL/2";
environment lab { require "site" == "bench"; }
tool readings { version "1"; mode deterministic; }
evidence series {
  tool readings;
  kind measurement_series;
  environment lab;
  max_age 3600;
  require "schema" == "EAL/typed-input/1";
}
reasoning rms_method {
  method "engineering/rms/1";
  rationale "Compute RMS deviation from zero over the supplied pressure samples.";
}
claim bounded_rms {
  statement "RMS pressure relative to zero is at most 5 kPa for this sample.";
  environment lab;
  proposition {
    subject "pump-A";
    quantity "pressure";
    unit "kPa";
    scope "sample-episode-v1";
    valid_from "2026-09-23T10:00:00Z";
    valid_until "2026-09-23T11:00:00Z";
    query {"origin":0};
    result "rms" <= 5;
  }
}
argument sampled_rms {
  conclusion bounded_rms;
  reasoning rms_method;
  evidence series;
  binding series;
}
\end{lstlisting}

The diagnostic's related vibration task supplies samples \(0.3\) and \(0.4\ \mathrm{m/s}\). With origin zero, the computed RMS is \(\sqrt{(0.3^2+0.4^2)/2}\simeq0.35355\ \mathrm{m/s}\), below the \(0.36\ \mathrm{m/s}\) claim limit. A counterexample retains those samples but records an origin of \(0.1\ \mathrm{m/s}\). It yields \(\sqrt{((0.3-0.1)^2+(0.4-0.1)^2)/2}\simeq0.25495\ \mathrm{m/s}\). That smaller result answers a different query. The checked claim is \texttt{unsupported} for that supplied observation, even though either scalar passes the numerical threshold \citep{earlTask}. The interpreter's result is conditional on the authored question, observation identity and method implementation; it does not certify that the physical measurement or model is adequate.

\subsection{Support, objections and bounded evaluation}
Let \(E(a)\), \(A(a)\) and \(P(a)\) be an argument's evidence, assumptions and premise claims. Local usability checks declared context and time, evidence acquisition identity and freshness, required predicates, and the method's typed input and output. The authored argument asserts a conditional relation between those checks and its conclusion. It may have alternative derivations; an observation reused by several arguments retains one identity. The evaluator rejects cyclic claim-premise dependencies and bounds programme size and graph relationships \citep{earlArgument}.

Objections form a support--attack graph, including objections against objections. Write \(U(n)\) for a node's local usability, \(P(n)\) for its required claims, \(A(n)\) for attackers, and \(D(c)\) for arguments deriving claim \(c\). Starting with undecided labels, the evaluator monotonically adds justified accepted or rejected labels to its least-information fixed point:
\begin{align}
\operatorname{Acc}(n)&\Leftarrow U(n)\land
  \bigwedge_{c\in P(n)}\operatorname{Acc}(c)\land
  \bigwedge_{a\in A(n)}\operatorname{Rej}(a),\\
\operatorname{Rej}(n)&\Leftarrow\neg U(n)\lor
  \bigvee_{c\in P(n)}\operatorname{Rej}(c)\lor
  \bigvee_{a\in A(n)}\operatorname{Acc}(a),\\
\operatorname{Acc}(c)&\Leftarrow\bigvee_{n\in D(c)}\operatorname{Acc}(n),
\qquad
\operatorname{Rej}(c)\Leftarrow\bigwedge_{n\in D(c)}\operatorname{Rej}(n).
\label{eq:labels}
\end{align}
The claim-premise graph is acyclic, while attacks and objection-support relationships may cycle. This is a defined conjunction/disjunction extension whose attack-only restriction agrees with grounded labelling for a finite Dung graph \citep{dung1995,earlArgument}. It does not generate objections from natural-language contradiction or rank competing methods.

For a usable argument, acceptance gives \texttt{supported}; a rejected or undecided derivation with no uncontested alternative gives \texttt{contested}. With no usable derivation the claim is \texttt{unsupported}; a failed environment check gives \texttt{out\_of\_scope}. These labels describe the supplied argument calculus. In particular, \texttt{unsupported} does not prove a proposition false. Method execution failure, an unmatched query and an unavailable observation remain distinguishable in the assessment record. Unit, time and scope checks reduce accidental reuse of a computation; engineers still choose the relevant quantities, assumptions and observations.

\subsection{Claim-scoped host and reuse boundary}\label{sec:host}
The later implementation moves an already authored argument into a host-owned catalogue. An operator pins source and method digests, context, assessment time and permitted claims. A trusted launcher authenticates the recipient's principal, validates finite exact claim identifiers, snapshots its grant sets and creates a separate MCP server exposing only assessment, explanation and completion for granted artefact--claim pairs. Recipient prose cannot introduce an ungranted source or claim; a pinned host invocation fixes the exact pair. A text-only model can request a schema-validated operation through the application host; the host, rather than the model, initiates MCP and executes the registered tools. The result of \texttt{finish} is recomputed against the saved assessment and remains the application status even if recipient prose contradicts it \citep{earlHost}.

The bounded claim packet (schema \path{eal2-claim-packet/2}) returns one status, source and method identities, context fingerprint, assessment and collection identifiers, and a short account of the selected argument, failed dependency or active objection. It limits the displayed arguments and objections and marks truncation; an authorised \texttt{explain} call retrieves a larger claim-specific direct trace. That trace identifies premise claims but does not expand all transitive premise evidence. A packet is at most 3,072 UTF-8 bytes and the explanation at most 65,536. Those limits constrain transport, not the number of semantic checks performed: the current \texttt{assess\_claim} path collects for the whole artefact before selecting one claim. An assessment ID alone is not an access grant. The recipient MCP route stores a principal-bound issuance record that survives server restart; the in-process task-family route instead retains issued IDs per instance and requires reassessment after its restart \citep{earlHost,earlFamilyHost}. A saved \texttt{finish} recovers a historically checked status; a claim about the current cluster needs fresh collection and reassessment. Source and observation digests detect a changed represented record; the packet explicitly labels evidence integrity \texttt{consistency\_checked\_not\_authenticated}. Physical sensor identity, collector integrity and truth of an authored warrant require separate evidence.

A host-owned family catalogue groups finite, typed parameter tuples, each mapped to an exact pinned artefact and selected claims. Its loader checks that each parameter has a distinct context path and a corresponding equality constraint in the selected claim's EAL environment. A recorded reviewer label and contract digest bind the declared review to the source, method, context, time and claims; the software cannot authenticate the reviewer or show that the review was independent. An unlisted tuple cannot inherit an existing verdict. Candidate retrieval ranks authorised families by lexical overlap with task text. It returns names and matched terms, without choosing parameters, granting access or asserting applicability. A trusted application then chooses explicit bindings and a claim, and \texttt{TaskFamilyHost} checks the family, artefact and claim grants before assessment \citep{earlFamilies,earlFamilyHost}. This avoids treating semantic similarity as an inference rule, at the cost of reviewing a new tuple when scope changes.

The architecture separates three possible errors. A source can misstate the original engineering question; an observation can be stale, counterfeit or irrelevant to that source; and a recipient can misreport an otherwise checked result. Host execution addresses specified checks in the second and third stages. It cannot infer omitted resource contention, model inadequacy or missing defeaters. Even an authorised, exact family can address a different question from a vague task: a supported status for a wrongly selected family is false support relative to that task, without a grant bypass. Applicability of a selected family to the user's intent needs trusted review. A deployed service additionally needs an authenticated endpoint that binds the principal and grant; the present trusted launcher and in-process route define that boundary but do not demonstrate arbitrary remote identity integration. Family candidate retrieval and exact binding are application APIs and are not yet exposed through the recipient MCP endpoint. Similarly, existing targeted objections and finite fixed-point labels do not amount to ASPIC+ construction, contrariness and rule preferences. If an engineering task requires those operations, a separately specified finite argumentation profile could be supplied through a versioned host method without changing EAL/2 syntax. It would need typed profile inputs, explicit construction and defeat semantics, a trace and binding from generated results to EAL claims, and independent checks before any claim of conformance \citep{modgil2014,earlArgument}.

\begin{table}[!htbp]
\centering
\small
\setlength{\tabcolsep}{4pt}
\caption{Claim delivery stages and the decisions each stage is permitted to make. Lexical nomination cannot authorise execution; recipient prose cannot replace a host status.}
\label{tab:host-architecture}
\begin{tabular}{@{}>{\raggedright\arraybackslash}p{2.5cm}>{\raggedright\arraybackslash}p{3.9cm}>{\raggedright\arraybackslash}p{4.7cm}@{}}
\toprule
Stage & Supplied material & Authority and output\\
\midrule
Nominate & Task text, permitted family metadata & Lexical suggestions; no binding or verdict\\
Resolve & Explicit family, typed tuple, claim and caller grant & Trusted application checks exact reviewed case and grants; pins artefact\\
Assess & Pinned source, context, method, observations and time & Host checks represented evidence and emits a claim status, IDs and short reason\\
Explain and finish & Principal-bound assessment ID and claim & Authorised direct trace; saved host status remains final\\
Communicate & Checked packet and task text & Recipient prose is retained separately and scored for fidelity\\
\bottomrule
\end{tabular}
\end{table}

\subsection{Resource contention as an inferential change}\label{sec:kubernetes}
An additional synthetic EAL/2 example represents a checkout-api deployment in the payments namespace on a shared Kubernetes node. It links a declared image digest and attestation record to a reported rollout, a five-minute \texttt{checkout-create} trial with a configured 200-requests-per-second target, 10,000 reported sampled outcomes, four sampled p99 latency values and a snapshot of requested CPU headroom. A Wilson lower bound and a finite temporal predicate check specified numerical conditions. The bounded canary claim concerns the recorded trial. A separate operating-basis claim proposes considering the same build for the next fifteen minutes under unchanged traffic and resource conditions; that future performance is an authored, defeasible inference, not a consequence proved by the sampled result \citep{earlKubeExample}.

The first engineer's evidence must connect source commit, build attestation, image digest, rollout generation, ready endpoint identities and observed requests and responses to the \emph{same} service episode. An assessor then needs coverage and units for the latency series, a defended sampling model for the Wilson calculation, and a trusted load-generator ledger of sent, eligible and selected counts, timing, achieved rate and sampling rule, as well as actual usage and throttling. The submitted target and 10,000 sampled outcomes alone establish neither achieved load nor unbiased selection. Kubernetes CPU requests affect placement, whereas CPU limits are enforced through throttling; a request-headroom snapshot does not bound a co-tenant's actual consumption \citep{k8sResources}. After a second engineer deploys a batch job, a record of the batch Pod and target sharing a node, with time-aligned CPU usage, target throttling, queueing and latency, supplies a specific challenge to the operating inference. A controlled batch-admission or isolation trial under matched traffic would strengthen causal attribution. A current cluster query alone cannot prove future p99 or exclude unobserved interference.

For the example's three authored revisions, the operating-basis status changes from \texttt{supported} to \texttt{contested} when the new batch observation activates an objection against the reasoning step, then to \texttt{supported} with a fresh trial and complete monitor nonfinding in a later record. That sequence alone does not show that the stated mitigation caused the change. The historical canary remains a bounded report of its earlier sample; its episode is identified by retained assessment inputs and time, not the recurring claim ID alone. Dropping the contention record makes the raw evaluator support the operating basis because no attack is submitted. A fixture-specific application gate retains that raw label but returns unresolved. Generic \texttt{assess\_claim} does not invoke this gate, so a production caller must attach an equivalent mandatory-acquisition policy before treating the raw status as current assurance. Before acceptance it checks the mandatory records' source, request, time and matching manifest trial, Pod and node identifiers; missing, stale, failed or incoherent monitor records fail closed. The response and p99 record envelopes now carry trial, Pod/node, window and series identifiers plus a payload digest; the gate recomputes and matches these against the manifests. Those producer-declared bindings check record coherence but do not authenticate the collector, prove physical-series origin or establish that the query observed all interference. An operational deployment needs separately verified acquisition coverage and event-triggered reassessment when the work mix or resources change \citep{earlKubeExample,k8sEviction}. This counterexample concerns evidence completeness and the source-to-world relation. A more elaborate argumentation solver cannot infer an omitted event from silence.

\begin{table}[H]
\centering\small
\setlength{\tabcolsep}{4pt}
\caption{What could warrant the synthetic Kubernetes claims and what must trigger reassessment. The listed records would still need independently checked provenance and coverage.}
\label{tab:kube-evidence}
\begin{tabular}{@{}>{\raggedright\arraybackslash}p{2.4cm}>{\raggedright\arraybackslash}p{4.1cm}>{\raggedright\arraybackslash}p{4.6cm}@{}}
\toprule
Claim dependency & Required evidence & Change that defeats reuse\\
\midrule
Running build & Signed build provenance, image digest, rollout generation, Pod and EndpointSlice identities & New image, rollout, selector or endpoint membership requires new binding\\
Bounded trial & Trusted generator ledger of offered, eligible and selected request IDs, achieved rate, outcomes, histogram and sampling rule & A changed traffic mix or missing sample coverage makes the trial insufficient for that question\\
Resource premise & Node and Pod IDs, requests and limits, actual usage, throttling, queueing and complete time-aligned monitor query & Co-tenant admission, placement, HPA or node-pressure change requires recollection\\
Near-term inference & Defended traffic and resource stationarity, watched changes and a fresh matched trial; intervention for causal attribution & An unobserved event or expired coverage suspends a current operating assurance\\
\bottomrule
\end{tabular}
\end{table}

\section{Evaluation design}\label{sec:methods}
\subsection{Questions and materials}
The design claim is that EAL/2 can express and execute the specified engineering dependencies. Executable examples and regression checks exercise that claim within the implemented profile; their pass count does not establish broad representational completeness or human readability. The model questions are different. Does EAL syntax, compared with meaning-equivalent JSON, improve a pinned recipient's status answers? Can relevant raw observations correct a wrong proposal without damaging a correct qualified proposal? Does supplying an interpreter conclusion change the answer? A rival explanation of any syntax difference is a prompt-surface effect involving length, tokenisation, familiarity or output compliance, rather than a general reasoning gain \citep{sclar2024}.

We distinguish the earlier relay pilot from the later notation diagnostic. The relay programme used six previously exposed synthetic tasks, two repetitions, 27 model/transfer conditions and 324 attempted endpoints; 323 completed. In its declared primary contrast, the same GPT-4.1 tool-enabled producer supplied already correct answers on twelve blocks to a GPT-4.1 nano recipient. With tool evidence the recipient retained 12/12; with the producer's answer alone it retained 8/12. The descriptive task-weighted difference was 33.3 percentage points, with a task-cluster interval of 0.0--66.7 points and an inconclusive prespecified decision. There were no producer errors in those twelve blocks, so this contrast has no correction denominator \citep{earlRelay}.

The later experiment used five exposed synthetic tasks: a circular defence, an independently grounded defence, a registered RMS calculation, a wrong-origin RMS observation and a sampled negative finding. A source-repair task was excluded because its interface could not be meaningfully matched in this diagnostic. Each task had eighteen conditions and two repetitions, giving 180 scheduled and attempted model calls. The recipient was snapshot \texttt{gpt-4.1-nano-2025-04-14}, at temperature zero with a 2,048-token output limit. Each call used a fresh message context, no tools and no repair. EAL and JSON packets had equal parsed meaning, observations and common reference; their natural input lengths differed and the reference used EAL vocabulary. A correct or deliberately incorrect proposed answer was fixed for each task/repetition stratum. The crossed views were answer only, relevant raw observations, irrelevant observations for another assembly, and raw observations plus interpreter conclusions. A fifth view supplied raw observations without a proposed answer \citep{earlProtocol,earlNotation}.

The experimental unit for a notation comparison is a matched task/repetition pair. Two repetitions of the same exposed template are dependent and sometimes contain different wrong proposal labels; they are not ten independent task types. There was no protected unseen-task sample. Assignment, model settings, prompts and references were frozen before generation, with a pre-generation amendment to completion-status scoring and the available-information reference. After a known-cost output-limit failure at call 30, an operational continuation finished the original schedule without changing prompts or retrying that failure. Its revised stop handling was documented after the event \citep{earlNotes}. This departure and task exposure prevent confirmatory population inference.

\subsection{Outcomes and estimands}
The strict primary score admits a response only if its whole text is a JSON object with exactly \texttt{claims} and \texttt{basis}, all requested claim identifiers and allowed statuses, and a list-valued basis. Full-information correctness \(Y^F_{ik}(n,v,a)\) requires all statuses to match the original complete-observation reference for task \(i\), repetition \(k\), notation \(n\), view \(v\) and proposal \(a\). Available-information correctness \(Y^A\) compares the same response with the frozen interpreter reference for the records actually supplied. These references coincide for raw-only, relevant-raw and assessed views. For answer-only and irrelevant-record views the interpreter has no relevant observation, so asserting support can agree with the hidden full answer while being unjustified by supplied information.

The finite paired notation contrast is
\begin{equation}
\widehat{\Delta}_{N}=
\frac{1}{5}\sum_{i=1}^{5}\frac{1}{2}\sum_{k=1}^{2}
\left[Y^F_{ik}(\mathrm{EAL},\mathrm{raw},\varnothing)
-Y^F_{ik}(\mathrm{JSON},\mathrm{raw},\varnothing)\right].
\label{eq:notation}
\end{equation}
The corresponding raw-evidence contrast holds the wrong proposal, task and notation fixed and compares \(\mathrm{answer+raw}\) with \(\mathrm{answer}\). It is a response difference under this prompt, not a naturally occurring error-correction rate. Every failed or malformed attempt remains in the correctness denominator. A detected false-support event means that an admitted response marks at least one requested claim \texttt{supported} when the selected reference does not. Malformed text and provider errors are excluded from that event count, so its zero would not establish absence of false assertions in unadmitted text.

For an admitted claim map \(m\) and a reference \(o\), define
\(\operatorname{Matches}(m,o)=(m\neq\varnothing)\land
\bigwedge_{c\in\operatorname{Ids}(o)}(m[c]=o[c])\).
Define \(\operatorname{FalseSupport}(m,o)=(m\neq\varnothing)\land
\bigvee_{c\in\operatorname{Ids}(o)}
(m[c]=\texttt{supported}\land o[c]\neq\texttt{supported})\).
\autoref{alg:score-paired-trials} summarises the retained strict scorer and matched contrast procedure using these two exact tests. The implementation additionally verifies frozen schedule and record integrity. The algorithm distinguishes the two reference questions and never repairs malformed model text.
\begin{algorithm}[H]
  \caption{Strictly score retained trials against both references and pair notation outcomes}
  \label{alg:score-paired-trials}
  \begin{algorithmic}[1]
    \Require Frozen task/repetition schedule \(S\), attempted records \(R\),
      full references \(F\), view-specific references \(V\)
    \Ensure Correctness indicators \(Y^F,Y^A\), detected false support,
      complete-pair contrasts
    \ForAll{scheduled row \(s\in S\)}
      \State \(r\gets\Call{FindAttempt}{R,s}\)
      \State \(C_s\gets\Call{ClaimIds}{F_s}\)
      \State \(m\gets\Call{AdmitStrictClaimMap}{r,C_s}\)
      \ForAll{reference \(o\in\{F_s,V_s\}\)}
        \State \(Y^o_s\gets\Call{Matches}{m,o}\)
        \State \(H^o_s\gets\Call{FalseSupport}{m,o}\)
      \EndFor
    \EndFor
    \ForAll{matched task/repetition block \(b\)}
      \If{both raw-only notation attempts are present}
        \State \(d_b\gets
          Y^F_{b,\mathrm{EAL,raw}}-Y^F_{b,\mathrm{JSON,raw}}\)
      \EndIf
    \EndFor
    \State \Return \(Y^F,Y^A,H,\{d_b\}\)
  \end{algorithmic}
\end{algorithm}

All 180 calls were attempted once. Eighteen responses failed the exact schema and one hit the output limit; every response retained model identity and usage. The summed reported usage was 1,919,842 input tokens (1,837,824 cached) and 29,816 output tokens. Configured-rate cost was USD 0.0660738, an estimate rather than an invoice. The archive stores raw records, hashes, scoring tables and the continuation record \citep{earlNotation}.

\section{Synthetic bias-attribution study}\label{sec:bias-stage-a}
\subsection{Construct and scope}
The earlier experiments test interpretation, source authoring and delivery of represented engineering evidence. The Stage A study tests whether model dispositions match stipulated process records in synthetic engineering episodes. Human reliance remains a separate empirical question. Automation bias has been studied as an attentional and reliance phenomenon in interaction with decision aids \citep{parasuraman2010}. Cognitive forcing reduced overreliance in one human--AI task, with differences in subjective acceptance and need for cognition \citep{bucinca2021}. A positive test can be informative under a task's prior distribution \citep{klayman1987}, verification cost can explain reliance \citep{vasconcelos2023}, and rejection of AI advice can itself be unwarranted \citep{dietvorst2015}. These alternatives motivate the study's explicit technical and organisational rivals.

The protocol defines a bounded episode containing an authored EAL/2 argument and a time-ordered process record \citep{earlBiasProtocol}. A proposed bias mechanism requires an exposed cue, an opportunity to act otherwise, evidence of uptake, a predicted directional signature, temporal order, a discriminating contrast and examination of rivals. Its disposition ranges from \texttt{NOT-APPLICABLE} and \texttt{REJECTED} through \texttt{HYPOTHESISED}, \texttt{COMPATIBLE} and bounded episode or comparative support to \path{CAUSALLY-SUPPORTED}; \texttt{UNDECIDED} records an unresolved basis. In Stage A, the author supplies both the process record and the reference disposition. The resulting score measures performance on that constructed reference. Static validation checks the represented declarations; the EAL host did not acquire these records or assess a person's cognition.

\subsection{Executed design and amendment history}
Stage A used twelve author-created families, each with four variants: a process-supported case, a stronger technical or organisational rival, a missing prerequisite and an insufficient process record. Missing prerequisites rotated among cue exposure, choice opportunity and uptake. The 48 variants remain dependent within family. Each case received one direct critic call and two separately prompted mechanism and rival calls for each of three requested model identities. This produced 288 model--case--route assignments and 432 API calls. Calls ran serially in seeded shuffled order; independent analysts received the same episode without seeing each other's answers.

The requested and returned identities were \path{gpt-4.1-nano-2025-04-14}, \path{gpt-6-luna} and \path{gpt-6-sol}. The retained fields show no identity drift within a model entry. Nano used a dated snapshot; the other two names are aliases. The request omitted an explicit reasoning setting for Nano, selected \texttt{none} for Luna and \texttt{high} for Sol, and allowed at most 4,096 output tokens per call. Thus model identity, reasoning allocation and provider rates vary together. The direct route used one call per case and the independent route used two; Stage A does not isolate topology at equal call allocation.

The original zero-call plans exposed case class through identifiers and source structure. Pre-call amendment 0.1.1 removed public class labels, aligned the candidate-objection scaffold and made raw disposition scoring primary. Its first pilot call produced an invalid response and stopped the runner. After that answer had been inspected, amendment 0.1.2 allowed independent assignments to continue after terminal invalid outcomes, while retaining each failure without retry. Pending calls and a prospective budget breach still stopped execution. The amendment retained the cases, prompts, requested models, scoring rules and ceilings. The first pilot request identifier and request hash matched the continuation freeze, so its original events were imported unchanged. The remaining 71 pilot calls completed, followed by 360 further calls. The pilot is part of the full 432-call cohort. This outcome-informed execution amendment makes the run developmental \citep{earlBiasStageA}.

PR~13 proposed a separate zero-call source refresh, also labelled 0.1.2, on a later implementation revision. That unexecuted proposal is not the continuation used here. The reported run uses freeze schema \path{eal2-bias-freeze/2}, execution rule \path{continue_after_terminal_outcome/1}, the retained Stage A freeze digests and historical source revision \texttt{430945c}. Later implementation changes cause current-source material checks to reject those historical freezes. Report reproduction therefore uses the pinned historical material without regenerating the executed inputs.

\subsection{Scoring and dependence}
An admitted direct disposition requires a valid response; an admitted independent disposition requires both analyst responses to pass the response and identifier checks. Synthesis requires agreement on status, mechanism, defect, correction and rival. Disagreement produces \texttt{UNDECIDED}. It intersects each evidence-role citation set and retains an agreed defeater, avoiding a complete chain assembled from two different partial accounts. The primary \emph{fully warranted} score then requires the reference status, mechanism or permitted non-attribution, required citations, rival, defeater, defect and correction. Invalid or absent dispositions score zero and remain in the assignment denominator. In this report, ``raw'' means before the additional author-labelled evidence-role gate, after response validation.

The secondary gate uses private author-labelled process roles and can demote an unsupported attribution. It is an oracle-assisted diagnostic whose deployment would require independent validation of how roles are obtained. Supported-case recall counts the right mechanism with a supported status on the twelve positive cases; it can succeed even when the stricter full disposition fails. Decisive coverage counts admitted non-\texttt{UNDECIDED} dispositions over all 48 cases. Before identifier validation, separately retained parsed analyst answers permit an additional audit of strong false attributions and wrong corrections. Rationale-based accusations still require a separate semantic review; no blinded human review of the Stage A rationales is retained.

For model \(m\), the topology contrast is
\begin{equation}
\widehat{\Delta}_{m}=\frac{1}{12}\sum_{f=1}^{12}
\left[\frac{1}{4}\sum_{v=1}^{4}
\left(Y_{mfv,\mathrm{independent}}-Y_{mfv,\mathrm{direct}}\right)\right],
\end{equation}
where \(Y\) is the binary fully warranted score. The retained analyser resamples twelve complete family differences with replacement 5,000 times, using seed 240924, and reports a descriptive percentile interval. This preserves the four related variants within family. The interval describes resampling of the authored cases; the cases were not sampled independently from an engineering population.

\subsection{Retained Stage A results}
All 432 calls reached terminal outcomes: 160 passed response validation and 272 were invalid. The 72-call pilot contained 24 valid responses. Across the full cohort, 98 of 288 route assignments produced an admitted disposition, and five met the full raw reference: two Luna direct, one Sol direct and two Sol independent. Table~\ref{tab:bias-stage-a} gives the distinct call, admission, correctness and coverage denominators. The author-labelled gate retained the same five fully warranted dispositions \citep{earlBiasStageA}.

% Generated by verify_stage_a.py from reproduced retained reports.
\begin{table}[!htbp]
\centering
\small
\setlength{\tabcolsep}{3pt}
\caption{Stage A on 48 synthetic cases per model and route. Valid calls pass response checks; independent admission requires two valid calls. Warranted is the full raw reference score. Recall counts supported answers on twelve positives; decisive counts non-\texttt{UNDECIDED} dispositions. Failures remain in every assigned denominator.}
\label{tab:bias-stage-a}
\vspace{4pt}
\begin{tabular}{@{}llrrrrr@{}}
\toprule
Model & Route & Valid calls & Admitted & Warranted & Recall & Decisive\\
\midrule
Nano & Direct & 1/48 & 1/48 & 0/48 & 0/12 & 1/48\\
Nano & Independent & 4/96 & 2/48 & 0/48 & 0/12 & 0/48\\
Luna & Direct & 10/48 & 10/48 & 2/48 & 5/12 & 9/48\\
Luna & Independent & 12/96 & 0/48 & 0/48 & 0/12 & 0/48\\
Sol & Direct & 41/48 & 41/48 & 1/48 & 9/12 & 41/48\\
Sol & Independent & 92/96 & 44/48 & 2/48 & 5/12 & 33/48\\
\bottomrule
\end{tabular}
\end{table}


The structured diagnostics recorded 145 calls with an unknown or duplicate evidence identifier in the rival-evidence field and 63 with an unknown defeater. Another 89 responses failed the rationale contract; their retained rationales exceeded 300 characters. Diagnostic counts can overlap within a call. No transport or provider error was recorded, every call had usage for the configured cost calculation, and no pending call remained unresolved. These are failures of the requested response and reference contract on the exposed cases; the invalid responses do not reveal how accurately a model would answer under a different contract.

The raw and gated scores recorded zero strong false attributions after response validation. Each model--route cell contained 36 assigned negative cases. The retained parsed analyst answers give a different observation: Nano direct made eight strong false attributions among 34 parsed negative-case answers; its independent analysts made twelve among 71. Those twenty answers failed subsequent validation and did not enter admitted raw dispositions. Sixteen parsed Nano answers also selected a wrong correction, eight in each route. The report therefore supports a failure-filtering result at the response boundary, while leaving the semantic safety of rejected answers and short rationales unresolved. Zero admitted false attributions cannot establish safety when many assignments yield no usable disposition. Positive recall and decisive coverage in Table~\ref{tab:bias-stage-a} expose that limitation.

The family-mean independent-minus-direct differences in fully warranted score were \(0\) for Nano, \(-0.0417\) for Luna and \(0.0208\) for Sol. The descriptive 95\% bootstrap interval was \([0,0]\) for Nano. For Luna it was \([-0.1042,0]\); for Sol it was \([-0.0417,0.0833]\). Nano's degenerate interval follows from zero fully warranted scores in every resampled family; it provides no evidence of equivalence. Every interval includes zero, and the sparse successes, outcome-informed continuation and unequal call allocation prevent a general topology-benefit claim. The observed route differences and model-specific contract failures remain descriptive findings on these fixtures.

The ledgers record 819,810 input tokens, zero reported cached input tokens and 229,225 output tokens. Applying the frozen uncached rates gives US\$2.4331466 across all 432 calls, including the imported pilot event exactly once. This configured estimate is below the US\$25 cumulative ceiling; it is not an invoice or an all-in cost. Summed serial request latency is 3,686.6752\,s, excluding preparation, interruption and review time. Neither the numerical price difference nor model aliases establish a small-model advantage or general ranking. Both pilot and full reports reproduce byte for byte from the retained ledgers and pinned material \citep{earlBiasStageA}.

\subsection{Unexecuted follow-ups}
The original follow-up implementation under \path{bias-followups/} reuses the exposed 48 episodes for EAL/2 versus mechanically derived typed JSON and two general critics versus mechanism/rival critics. Its pilot and full plans contain 128 and 768 calls and have no retained provider outcomes. Equal call counts and output ceilings control assigned calls, while input length, reasoning use and billed tokens can still differ. Reusing the Stage A families makes any later outcome developmental.

The separate plans under \path{bias-next/} instead require fresh, mutually disjoint family sets for matched JSON and equal-call topology comparisons, with independent reference review and separate freezes. Prior benchmark findings on debate and ensembles vary with the task and design \citep{du2024,smit2024,ma2025}; they do not supply the missing Stage A control. A notation effect requires equivalent represented information. A topology effect requires an explicit allocation contrast, and an isolated model comparison requires comparable settings. These plans remain unexecuted and lack the required material receipts \citep{earlBiasNext}.

The human-decision protocols under \path{bias-human-trial/} and \path{bias-next/human-decisions/} remain distinct, unexecuted design records. Both motivate pre/post decisions, correct and incorrect advice, presentation controls and independent case review; their assistance arms and primary contrasts differ and must be reconciled before registration. The newer plan requires fixed reviewed agent packets prepared before recruitment and excludes live API calls during participant observation. Four synthetic administrative tests of the earlier trial tooling establish allocation and scoring behaviour only. No participant outcomes, reviewed human case manifest or ethics approval are retained. A prospective randomised human trial, with a justified precision target and frozen primary contrasts, is required to estimate changes in human decision quality or reliance errors.


\section{Results}\label{sec:results}
The earlier studies in this section and Section~\ref{sec:delivery} measure represented engineering status and recipient communication. Section~\ref{sec:bias-stage-a} reports the separate synthetic bias-attribution cohort. Human decision benefit remains unmeasured across these records.
Table~\ref{tab:conditions} retains the attempt denominator and separates the full and supplied-information references. The strict raw-only notation comparison yielded EAL 4/10 and JSON 5/10, or \(\widehat{\Delta}_{N}=-0.10\). Four matched pairs were correct in both formats, one only in JSON and five in neither. No task showed an EAL-only success; the single discordant pair arose in the sampled-negative-finding task. This describes these ten paired attempts and does not estimate a population JSON advantage.

\begin{table}[!htbp]
\centering
\small
\setlength{\tabcolsep}{4pt}
\caption{Strict complete correctness, with ten attempted calls per notation and view. Correct and incorrect describe the injected proposal against the full-information answer.}
\label{tab:conditions}
\begin{tabular}{@{}lrrrr@{}}
\toprule
View and proposal & EAL full & JSON full & EAL available & JSON available\\
\midrule
Raw only; no proposal & 4/10 & 5/10 & 4/10 & 5/10\\
Correct; answer only & 10/10 & 9/10 & 2/10 & 2/10\\
Correct; relevant raw & 2/10 & 7/10 & 2/10 & 7/10\\
Incorrect; answer only & 0/10 & 2/10 & 4/10 & 4/10\\
Incorrect; relevant raw & 2/10 & 5/10 & 2/10 & 5/10\\
Incorrect; assessed & 7/10 & 8/10 & 7/10 & 8/10\\
Correct; irrelevant & 6/10 & 6/10 & 0/10 & 0/10\\
Incorrect; irrelevant & 2/10 & 3/10 & 3/10 & 0/10\\
\bottomrule
\end{tabular}
\end{table}

With a deliberately wrong proposal, relevant raw records produced two EAL and five JSON complete corrections, compared with zero and two for answer-only. Supplying interpreter conclusions along with the raw records raised complete correctness to seven and eight. These differences involve supplied information as well as changes in prompt length, and cannot isolate independent calculation by the recipient. In the correct-proposal stratum, EAL answer-plus-raw scored only 2/10 versus 10/10 for answer only; JSON scored 7/10 versus 9/10. Under the supplied-information reference, answer-only correctness was 2/10 for a full-information-correct proposal in either notation, because the proposal itself supplied no observation. Describing that reference change as damage to correct empirical support would invert the evidence requirement.

The irrelevant-record control exposes a separate failure. Across its four representation/proposal arms, only 3/40 attempts were correct for information actually available, and 26/40 contained detected false support in admitted responses. None of the 28 proposals wrong under that reference was corrected. Values from another assembly could appear numerically plausible, but they were not observations for the requested scope. In raw-only trials both representations failed the circular-defence and wrong-origin tasks in both repetitions; each succeeded on the independently grounded defence and matching RMS task twice \citep{earlNotation}.

The prespecified strict output contract counts the eighteen malformed responses as failures. A later, explicitly exploratory readout accepts a strictly parsed raw JSON claim map with extra top-level fields and ignores the \texttt{basis} shape. It changes raw-only results to EAL 4/10 and JSON 6/10. It also increases detected false-support events across all conditions from 69/180 to 77/180. Fence-only normalisation changed no primary count. Relaxing output shape therefore cannot reveal an EAL advantage in this sample; it changes the measured endpoint and was chosen after schema failures had been seen \citep{earlNotation}.

\section{Host delivery and its present evidence}\label{sec:delivery}
\subsection{Earlier checked-packet recipient pilots}
Two subsequent developmental runs each made 48 one-shot requests to the same GPT-4.1 nano snapshot. Four selected synthetic roots had three related evidence states each. Every state was rendered as raw EAL, an EAL-derived JSON graph, EAL with an application instruction, or an earlier checked host packet. This packet predated the claim-scoped \texttt{/2} packet and the principal-bound family route in Section~\ref{sec:host}; neither new feature was exercised by those model calls. The independent brief-level reference and interpreter agreed for the twelve states, but no second masked source review or independently implemented equal checker was part of the live blocks \citep{earlArtifactPilot,earlArtifactReplication}.

\begin{table}[!htbp]
\centering
\small
\caption{Exact recipient status and format on twelve related states in each earlier 48-call developmental run. Counts are descriptive; the checked packet already supplied the host's answer.}
\label{tab:artifact-pilots}
\setlength{\tabcolsep}{4pt}
\begin{tabular}{@{}lrrrr@{}}
\toprule
Route & First exact & First false & Repeat exact & Repeat false\\
\midrule
Raw EAL & 3/12 & 2/4 & 6/12 & 2/4\\
Derived JSON & 8/12 & 3/4 & 8/12 & 3/4\\
Application instruction & 5/12 & 1/4 & 5/12 & 1/4\\
Checked packet & 10/12 & 0/4 & 10/12 & 0/4\\
\bottomrule
\end{tabular}
\end{table}

The earlier checked packet required 3,754 recipient input tokens over twelve calls, compared with 17,855 for raw EAL in each schedule. Those are reported provider tokens, not total workflow resources. The recipient still contradicted the host on two supported inspection states in each run; a correct application result therefore retains the host's status rather than adopting the recipient's field. The first packet also produced at least one correct label with an inaccurate causal explanation under subsequent qualitative inspection. No blind explanation score, acquisition cost, host latency or authoring effort was recorded. The JSON graph was generated from EAL source rather than independently authored, and some of its input tokens were cached. These runs concern four selected roots; they do not estimate a gain from the new host or EAL notation \citep{earlArtifactPilot,earlArtifactReplication}.

\subsection{New implementation checks and cross-model comparison}
Two developmental schedules crossed five routes with three configured model classes: raw source and records; the same input with a short skill instruction; a model-originated JSON assessment request through the host followed by a possible second call; the compact claim-scoped EAL packet; and the same final-answer contract using an independently implemented generic graph checker. The strong reasoning-model route is the raw condition on that model class. The request route validates a model-selected artefact and claim against a fixed allowed set; the correct candidate description was placed first in every route brief, so this tests menu/request syntax rather than family retrieval or new parameter bindings. The previously exposed four-root corpus supplied 180 first-call scenarios and a separate, singly authored three-root cohort supplied 135. All source, observation, prompt, checker and order identities were frozen before calls. An initial 180-case freeze stopped after three attempts (two completed and one returned-model identity mismatch); those attempts and their USD 0.0050428 configured-rate cost remain a separate diagnostic. Revised freezes allowed the observed dated nano identity and completed 180/180 and 135/135 scheduled cases without a retry or failed provider attempt \citep{earlCampaign,earlCampaignLive}.

Offline, the generic checker and EAL reference agreed on 12/12 scored states from the older pilot and 9/9 from the new cohort. Five newly authored challenge pairs passed in both systems. Sixty generated older-pilot and 45 new-cohort envelope mutations exercised the checker alone, so those mutations do not establish EAL/checker equivalence. The new cohort has no independent human reference, source or comparator review. The resulting parity establishes agreement on selected represented cases; it neither certifies the reference against real engineering systems nor measures model reasoning. For both checked routes, the final host status inherits the oracle-matched reference by construction. Thus every checked final status matches that reference and none finally falsely supports an adverse state, but this is an enforcement property of the runner. Recipient format and explicit agreement with that status are distinct outcomes \citep{earlCampaignLive}.

\begin{table}[!htbp]
\centering
\small
\setlength{\tabcolsep}{4pt}
\caption{Developmental cross-model calls. Raw and skill columns count strict exact recipient statuses. EAL-host and equal columns count strict recipient JSON with a status consistent with the host-owned answer, before explanation review; their final application status is oracle-matched by construction. Request counts are valid first-turn host requests, not valid final answers.}
\label{tab:cross-model}
\begin{tabular}{@{}llrrrrr@{}}
\toprule
Cohort & Model & Raw & Skill & EAL-host & Equal & Request\\
\midrule
Exposed, 12 states & Nano & 3/12 & 2/12 & 8/12 & 6/12 & 3/12\\
 & Luna & 10/12 & 10/12 & 8/12 & 11/12 & 12/12\\
 & Sol & 10/12 & 12/12 & 12/12 & 12/12 & 12/12\\
\addlinespace
New, 9 states & Nano & 1/9 & 1/9 & 0/9 & 3/9 & 5/9\\
 & Luna & 8/9 & 8/9 & 5/9 & 9/9 & 9/9\\
 & Sol & 7/9 & 7/9 & 9/9 & 9/9 & 9/9\\
\bottomrule
\end{tabular}
\end{table}

The exposed cohort had two false-support outputs in raw and two with the skill (one each for nano and Luna in each arm); the new cohort had none in those arms on its selected adverse states. For Luna, the EAL packet elicited a status contradicting the host in four of twelve valid pilot responses and four of nine valid new-cohort responses, whereas the equal-checker packet had one and zero contradictions. Nano returned malformed EAL responses in 4/12 and 8/9 cases, and malformed equal-checker responses in 6/12 and 6/9 cases. Sol gave valid, consistent checked responses throughout. These differences describe selected related states and unequal packet contents, not evidence of a notation-specific mechanism. They also show why the host-owned final status cannot be replaced by the recipient's prose \citep{earlCampaignLive}.

The request arm elicited 27/36 valid first-turn host requests in the exposed cohort and 23/27 in the new cohort. Its frozen second turn retained a system instruction demanding another operation request while the new user turn demanded final claims and explanation. All 63 final answers failed the strict final schema, including those after successful host calls; this is a prompt-conflict diagnostic, not a clean test of routed delivery. A corrected route was separately frozen and run. The 180-case run made 207 model calls: 257,339 input and 25,359 output tokens. The 135-case run made 158 calls: 177,085 input and 17,407 output tokens. Summed per-scenario API times were 1,923.205 and 1,448.013~s. Median per-scenario latency by model ranged from 7.68 to 11.82~s and from 7.92 to 11.59~s, respectively, including routed second turns. At configured rates, model charges were USD 0.3798973 and USD 0.2558221, or USD 0.6357194 combined. These are provider usage and call timings, not invoices, wall time, cluster-acquisition cost or all-in cost \citep{earlCampaignLive}.

Two later, separate route-only freezes corrected the second-turn instruction while retaining the same fixed first-turn menu and synthetic cases. Across 36 exposed-root and 27 new-root attempts, 51/63 requests produced an accepted host assessment and a strict recipient answer whose status agreed with the host: nano 9/21, Luna 21/21 and Sol 21/21. The twelve other nano first turns were rejected as nine incorrect or unauthorised requests and three malformed requests, not safe negative decisions. All 51 accepted statuses matched the singly authored reference, including seventeen adverse states, because the host supplied the checked answer; these are recipient-consistency counts before explanation review. The two revised runs made 114 model calls, with 83,911 input and 5,866 output tokens, zero provider failures or retries, and USD 0.0976742 in configured-rate model charges. Median per-case API times by nano, Luna and Sol were 11.179, 18.507 and 18.589~s. The first-turn samples are separate draws, so the change from the original 63 malformed final answers diagnoses the prompt procedure rather than a learned recipient improvement. It still tests a prefixed fixed menu, not candidate retrieval, typed family selection or native MCP tool use \citep{earlRouteV2}.

An exploratory post-call review sampled one state from each of seven roots across fifteen model--route conditions (105 outputs from the original flawed route schedule and four other arms). Two separate de-identified AI review passes used the same single-author dossiers. The adjudicated categories were 49/105 faithful, 13 incomplete, four materially false and 39 malformed or absent; reviewers agreed on the faithful binary for all 105 and differed on one non-faithful category. Judged faithful counts by route were raw 16/21, skill 15/21, EAL-host 4/21, equal checker 14/21 and the flawed request route 0/21. The subset rule was documented after some outputs were observed, and response text could reveal an arm despite withheld linkage. The reviewers saw a full task dossier rather than each recipient's exact exposed input. A cautious statement about missing input could therefore be penalised despite reflecting the recipient's actual view. The compact packet also omits some transitive premise evidence. These are exploratory end-to-end explanation-completeness signals on dependent synthetic states, not human adjudication, a model-fault attribution, or a complete 315-case fidelity rate. They cannot supply the denominator for all-in cost per correct accepted decision \citep{earlExplanationReview}.

A separate exploratory review assessed all 63 corrected route finals against a rubric fixed before those outputs were scored. Its packet retained the actually delivered host packet alongside author source and observations, while withholding model linkage. Two independent AI reviewers agreed on 59/63 exact categories and 61/63 faithful-binary labels; a third agent resolved four differences to 11 faithful, 38 incomplete but nonmisleading, two materially false and twelve malformed or absent. Thus all 51 accepted recipient statuses matched the host, yet only 11/63 route attempts (11/51 accepted answers) met this strict evidence-level explanation rubric. Generic claims that the host accepted a route often omitted a decisive record, objection or calculation. The packet itself can omit transitive premise detail, so this is an end-to-end completeness finding on authored synthetic cases, not a causal diagnosis of recipient reasoning or human-validated fidelity. The agent review time omits source authoring, independent verification and live acquisition and cannot establish all-in cost \citep{earlRouteV2Explanation}.

A separate synthetic Kubernetes host/retrieval study specified seven evidence revisions before its local timing run. It separates a historical test-record claim from a provisional operating-basis claim. The host matched the author's seven expected operating statuses and decisive causes and seven expected historical-record statuses. A fresh pressure report contested only the operating basis; stale load or monitoring coverage left its earlier record supported but made the operating basis unsupported. A separate reservation report answered the named pressure objection. This corrected case supersedes an earlier fixture in which an attack on a literal report claim conflated a later operational challenge with falsification of the past record. The mean compact packet occupied 1,553 UTF-8 bytes versus 6,042 for its direct trace, a mean ratio of 0.257. In twenty local repeats per revision, median combined assessment, explanation and completion time ranged from 161.0 to 174.4~ms. These timings include local Python and filesystem work, not MCP network transport, live cluster collection or model latency \citep{earlKubeOffline}.

Ten singly authored candidate queries gave lexical retrieval 8 true positives, 6 false positives and 1 false negative at a limit of three: micro precision 0.5714 and recall 0.8889. Top-one accuracy was 7/8 among queries with a target; one of two no-target queries produced no suggestion. The broad Kubernetes query returned an unwanted candidate, and a traffic-switch synonym missed its intended family. Exact resolution rejected the unlisted \texttt{prod\_west} tuple and an unauthorised family--claim pair despite candidate suggestions. These are developmental labels from the fixture author, not independent retrieval relevance judgements. Packet size is not token usage, and 7/7 status agreement tests consistency with an authored graph rather than whether the authored inference or a cluster observation is correct \citep{earlKubeOffline}.

A distinct developmental selection schedule presented lexical suggestions to three model classes on twelve synthetic tasks, yielding 36/36 completed one-shot calls. The model saw task text, trusted scope and candidates, but no EAL source, evidence revision or checked answer; it selected a family, typed cluster/namespace binding and claim, or abstained. The trusted exact resolver owned the subsequent status. Selection or appropriate abstention was correct for nano on 9/12 tasks, Luna on 10/12 and Sol on 10/12. All three selected the same latency family across baseline, contention and stale-coverage evidence revisions; the host returned \texttt{supported}, \texttt{contested} and \texttt{unsupported}. All three unlisted \texttt{prod\_west} tuple requests were denied. A missing lexical candidate for a traffic-switch synonym led all three to abstain on a legitimate task. On two ambiguous or broad tasks nano selected an authorised but irrelevant latency family, and its checked \texttt{supported} result gave two task-level false supports. This was an applicability failure rather than a grant bypass; Luna and Sol abstained on those tasks \citep{earlFamilySelection}.

The corrected frozen input digest for this selection study is \path{b5ef86e16f16d539d75d0632497b73862a6f6ed4c1ea9cd42d6fdb313244c7cc}; an earlier audited freeze was superseded with zero paid calls. The retained archive records 10,938 input and 2,326 output tokens, no retries or failed calls, and USD 0.0195864 in configured-rate model charges. Median per-call API times were 7.339, 8.316 and 8.899~s for nano, Luna and Sol. Its exact freeze and usage ledger were independently audited, but task relevance, source and statuses were labelled by one fixture author. These results demonstrate selected text-only candidate handling and exact denial on synthetic records; they do not establish field selection accuracy, live cluster truth, native MCP model use or faithful explanation \citep{earlFamilySelection}.

The separate original 96-call negative-revision protocol specifies 48 one-shot interpretation calls and 48 linked drafting/revision calls. The original 960-call mechanism protocol specifies 24 roots, ten content arms, two representations and two reference renderings, plus 192 deterministic executions over four record-integrity variants in two systems. These are different exposures from a principal-bound, host-owned final answer. The original 800-call deployment protocol says its base count excludes an optional candidate grammar comparison. Its stated fixed-source subtotal of 384 is arithmetically inconsistent with the written full factorial: eight roots, three states, two formats, two instruction settings and eight model--delivery cells require 768, which would make 1,184 calls after its 96 author and 320 recipient calls. An 800-call execution therefore requires a versioned balanced subset of fixed-source cells, with some interactions unavailable, or a changed total. The existing typed method and targeted objections express the finite negative-finding tasks without a grammar change. New schedules must identify their amendments and input freezes; historical frozen calls remain historical evidence \citep{earlProposals,earlAmendment}.

\subsection{Three amended developmental schedules}
We retained the original protocol files at version 0.1.0 and separately versioned three synthetic schedules. Their reference claims concern the submitted finite records, not authenticated physical conditions. The sources and records came from selected author-created families. The pre-call gates differed: the 96 schedule had a second AI rule review; the 960 schedule had a masked AI label review and a separate prompt review; the 800 schedule had two unmasked AI label reviews. None supplied the original requirement for two masked human adjudicators or an independent human brief-to-source validation. Thus the observations below are developmental; repeated calls within one root are dependent, and matching EAL and comparator labels on the frozen fixtures is a gate rather than an empirical language advantage \citep{earlThreeStudyMethods}.

The negative-revision amendment used eight briefs from one author and a generated scaffold. They share the registered sampled-negative method but vary which positive route survived, which objection was answered and whether a covered absence could itself support the bounded property. A later substitute conveyor record bears an author-asserted correction sequence and digest of the earlier record; the acquisition gate checks that assertion and timestamp, not the physical history. EAL and an independently coded finite JSON graph checker matched the authored three-claim reference on all 24 canonical states and on pre-call acquisition and method mutations. Stage A assigned a fresh nano call to each root, state and representation (48 calls). Stage B assigned sixteen linked three-turn authoring threads, one per root and representation (48 calls). Every prompt, response, model identity and usage was retained \citep{earlNeg96}.

All 96 negative-revision calls completed without a provider failure or retry. In Stage A, raw EAL yielded the exact three-claim vector in 10/24 cases and the generic rendering in 3/24, while neither completed an entire root's three-state sequence. On the adverse state both views were exact in 0/8 roots and each had six of eight responses with false support; across all states those counts were 9/24 EAL and 7/24 generic. The EAL prompt was longer (median 6,310.5 versus 3,698 canonical-message UTF-8 bytes), and content and rendering changed together. In Stage B, the frozen strict checker accepted none of the 24 EAL or 24 generic submissions. The EAL outputs omitted a required header or used invalid tokens; the generic outputs were JSON but supplied route objects where reader-ID strings were required. The authoring prompt gave no complete executable exemplar, and the generic checker returned an unhelpful type error, so this stage diagnoses an authoring-interface failure rather than relative language capacity. The run reported 174,260 input tokens, including 52,608 cached tokens, 31,795 output tokens, US\$0.0261984 at configured provider rates and 10.89\,s median API time per call. These charges omit source and review effort \citep{earlNeg96}.

A mechanism rubric was frozen after the campaign began but before either reviewer opened its responses. Two AI reviewers then scored all 48 Stage A responses using each actual exposed prompt. They agreed on 47 strictly parsed answers, thirteen exact three-claim vectors and the same single response whose prose faithfully described sampled tests and a complete non-detection limited to its scope without claiming a violation. Hence only 1/48 was both exactly labelled and faithful under this rubric. Their subcriterion judgments differed in ten cases; one reviewer had earlier used a different categorical rubric and learned the other review's aggregate count before its comparable pass. These are exploratory agent ratings, not blinded human fidelity or a population estimate \citep{earlNeg96Explanation}.

The mechanism amendment crosses 24 selected roots in four task families with ten supplied-information arms, two meaning-equivalent EAL/parse-derived JSON representations and two reference renderings, giving 960 fresh one-shot nano assignments. The arms distinguish admissible and invalid raw records, a wrong proposal, intermediate eligibility or method information, an actual checked status and an identity-mismatched earlier verdict. A wrong \texttt{out\_of\_scope} proposal is conspicuous for these in-scope tasks; rejecting it is a weak test of resistance to plausible false support. The longer reference includes the compact core and additional documentation, making the contrast a whole-prompt intervention rather than a pure length effect. Before calls, 24 roots crossed with four input or gate variants and two systems gave 192 deterministic executions with matching authored references. The second system is a separately coded but same-author finite brief-rule checker, and the JSON rendering comes from the EAL parse; these dependencies limit any notation-specific inference. One AI reviewer derived the valid and invalid labels without the author's labels or model outputs, and another checked the frozen prompts and gates \citep{earlThreeStudyMethods}.

All 960 mechanism assignments were attempted across the original segment and three documented continuations: 942 strict completions, four malformed answers and fourteen provider failures. Eleven response-less failed requests have unknown actual charges. Under the reference for information available in each arm, 460/960 attempts gave the exact status and 130/960 admitted a false \texttt{supported} status; malformed and failed responses were not treated as safe abstentions. Valid raw records without a proposal gave 78/96 exact and four false supports, while invalid raw records gave 17/96 exact and 65 false supports. With valid raw records, a correct proposal gave 86/96 exact versus 36/96 for the conspicuously wrong proposal; supplying the checked status with that wrong proposal gave 76/96. An invalid record plus an identity-mismatched earlier status gave 18/96 exact and eighteen false supports. Across all arms EAL had 237/480 exact statuses and 60 false supports, versus 223/480 and 70 for parse-derived JSON. The complete selected-root paired mean differences were +0.0292 for EAL versus JSON and +0.0875 for compact versus longer, content-changing reference; neither identifies a syntax or context-length effect. The campaign recorded 6,008,007 input tokens (4,812,032 cached), 110,400 output tokens, a median 11.07\,s API time and US\$0.2840583 known configured-rate study-call charges, plus US\$0.0000069 in three separately reported probes. Its conservative reserve for response-less charges was US\$0.0301785, not an observed bill. The one-author synthetic fixture family, weak wrong-label test and absent human explanation adjudication bar a population or faithful-decision cost claim \citep{earlMechanisms960}.

The deployment amendment assigns 800 decision slots by selecting two complementary representation-by-instruction cells in each of 24 root/state blocks, counterbalanced across blocks. Stage~1 schedules 384 fixed-source slots over EAL and a separately checked graph, direct interpretation, a preassessed host packet and, for tool-capable models, a native Chat Completions function request to a validating host. The function condition records selection of a pinned artefact and claim and uses the host's final status; its successful attempt uses one model round trip, with any declared provider retry counted and charged separately. It does not test a second-turn model explanation or a live MCP service. All Luna and Sol cells in this amended deployment schedule use reasoning effort \texttt{none}, including native-function cells, after a separate Sol high-effort tool probe returned HTTP~400; the earlier text-only cross-model cohort used Sol at high effort. Stage~2 schedules 96 linked author and revision slots across eight roots, two formats and two author models, followed by ten isolated recipients of each of 32 final author products (320 slots). The eight briefs vary domain contracts but share a two-positive-route structure: an alert targets the primary route, while an independent alternative or exact answer can retain support. Two unmasked AI reviewers accepted the 24 finite synthetic labels; EAL and the graph checker matched those labels and passed 120 paired acquisition mutations. The half-fraction cannot identify the full format-by-instruction interaction within a root/state block \citep{earlThreeStudyMethods}.

All 800 assigned model-call slots in the fresh deployment schedule completed in 804 actual API requests. Four response-less HTTP~500 first attempts, all in Stage~1, were followed by the frozen identical-payload single retry; none became a failed slot. Table~\ref{tab:deployment-routes} reports the fixed-source accepted-status outcome, not explanation quality. The host-owned and native routes use the checked status by construction, whereas direct answers are model-produced: 50/144 exact overall, comprising nano 2/48, Luna 2/48 and Sol 46/48 (23/24 in each representation). All 96 native requests selected the pinned operation and claim. No strictly admitted false support was scored, but malformed or unaccepted text is not proof of safe abstention. The native result supplies no recipient-authored explanation, and the route-specific model charges omit host and acquisition costs \citep{earlDeployment800}.

\begin{table}[!htbp]
\centering\small
\caption{Amended 800-slot study, fixed-source Stage~1. Charges use configured model rates; route median is API time per completed assigned slot, excluding source authoring and host execution.}
\label{tab:deployment-routes}
\begin{tabular}{@{}lrrrr@{}}
\toprule
Delivery & Assigned & Exact status & Model US\$ & Median s\\
\midrule
Direct interpretation & 144 & 50 & 0.4506870 & 9.794\\
Preassessed host & 144 & 144 & 0.1169620 & 9.782\\
Native function and host & 96 & 96 & 0.0633402 & 9.226\\
\bottomrule
\end{tabular}
\end{table}

The 96 linked author turns completed across 32 groups; every second candidate changed from its first, and 21 third candidates changed from their second. Both independent AI source reviews judged all 32 final products invalid against the exposed executable contract: the EAL outputs lacked a legal language header and the graph outputs used an object where the checker required a claim identifier. The prompts described constructs but supplied no complete executable exemplar. The frozen host therefore marked all 320 actual recipient packets \texttt{unavailable}. Ninety-three recipient texts met the JSON output shape and 227 did not, but none had an accepted status; the 0/320 accepted count is not 320 incorrect engineering labels. Across all 800 slots the study reported 968,533 input tokens (45,808 cached), 112,900 output tokens, US\$1.16650454 known configured model charges and a US\$0.1456512 conservative reserve for possible charges on the four response-less first requests. The author and recipient model calls accounted for US\$0.42032954 and US\$0.11518580 of the known total. Median assigned-slot API times were 15.726\,s for author turns and 8.925\,s for recipients. These are performance under the frozen prompts, not evidence that the underlying engineering cases lack a valid EAL/2 representation \citep{earlDeployment800}.

For a preselected 48-response explanation subset, the retained packet included each actual recipient prompt, host packet, answer and author-specified reference. Two separately frozen AI reviews agreed on the overall ratings for all 48: four faithful responses and seventeen shape-valid answers. The four faithful answers were all fixed direct Sol responses among sixteen raw-source cases. All sixteen checked-host statuses matched the synthetic reference, but neither reviewer rated any of their recipient explanations faithful; the other sixteen selected cases received an unavailable authored-source packet. The reviewers disagreed on some component criteria, including qualification in fifteen cases, without forced adjudication. This is an exploratory, unmasked AI fidelity assessment under a rubric frozen during the campaign before reviewer exposure, not an independent human quality estimate or an effect estimate between routes. In particular, a host-correct final status does not make the recipient's explanation faithful \citep{earlDeployment800Explanation}.

In the developmental route comparisons, an adverse event is an admitted \texttt{supported} output when the author-specified reference requires another status. A confirmatory endpoint needs an independently reviewed reference. The full decision endpoint requires a task-relevant authorised family, tuple and claim binding, adequate evidence scope, the reference-correct status, the output contract and a separately adjudicated faithful explanation; malformed answers, failures and retries remain in attempted denominators. Let \(D_i=1\) only if decision \(i\) meets all these conditions against an independently reviewed reference, and zero otherwise. For \(N\) evaluated decisions over the declared reuse horizon, the cost per correctly accepted decision is
\begin{equation}
 \frac{C_{\mathrm{initial}}+C_{\mathrm{calls}}+C_{\mathrm{operations}}+C_{\mathrm{review}}}
 {\sum_{i=1}^{N}D_i}.
 \label{eq:total-cost}
\end{equation}
The initial term includes human authoring and source review for every new source, scope and evidence-policy revision, measured in time at a declared valuation. The calls term includes all attempted author, AI source-review and recipient model calls, retries and provider failures. Operations include evidence acquisition, host/checker compute and transport. Review includes recipient-input-matched explanation adjudication; research-only review is itemised separately from any operational review a deployed workflow requires. Initial costs are included once over that horizon rather than omitted from a per-call estimate; a zero denominator leaves the ratio undefined. The reported schedules have neither complete human effort nor human explanation adjudication, and the configured short-context token rates are not invoices and leave response-less failures' actual charges unknown \citep{openaiPricing2026}. The analysis therefore leaves total cost per correct accepted decision unavailable while reporting observed study-call provider charges separately \citep{earlCampaign}.

\section{Interpretation and limitations}\label{sec:discussion}
The implemented checks have a clear use: an engineer can ask which observation, query, scope, assumption or challenge caused a status and can reassess after a recorded change. Typed method contracts prevent a numerically passing result for the wrong formal query from silently supporting the claim. That is a software behaviour claim supported by the specified semantics and repository checks. Whether the authored relation from a measurement to a physical conclusion is sound still requires domain evidence. Reusing the same observation through two arguments does not create two independent measurements; the calculus does not assign numerical credibility to agreement among arguments.

What can the model trials establish about the notation itself? In this diagnostic they establish only a finite paired outcome for one model on five exposed tasks. Meaning-equivalent EAL and JSON differ in tokenisation and length, and the common reference uses EAL terms. Even a larger observed difference in this design would combine those surface mechanisms. Prompt-format sensitivity elsewhere supplies a reason to compare plausible formats; it does not predict EAL's direction of effect \citep{sclar2024}. The present strict and exploratory comparisons give no positive EAL notation result. A fully structured JSON interchange remains a credible design alternative for the same typed representation.

Computed interpreter conclusions are different from raw measurements. Their larger correction counts show that the tested recipient can often use a supplied status; they do not show that it independently applied the method or checked the scope. The earlier relay benefit likewise concerns preservation of a correct producer's qualified answer. The checked-packet pilots and new cross-model responses show that communication can fail after a correct host assessment. The frozen request-route conflict also demonstrates an application-prompt failure after successful host calls. These observations separate preservation, recovery from injected wrong labels, host execution and recipient communication. A host-owned final result enforces the last boundary for represented statuses, but the engineering warrant and acquisition chain remain independently reviewable.

The amended schedules expose different bottlenecks. The 960 one-shot information interventions left 130 admitted false supports and a large valid-versus-invalid raw-record gap; model prose cannot be treated as an admissibility gate. In the 800 fixed-source stage, the preassessed and native routes retained all checked final statuses because the application owned the decision, while direct Sol gave 46/48 exact responses under this prompt and nano and Luna each gave 2/48. The authored stage then failed the executable source contract in every final candidate, so its 320 recipient packets had no valid assessed claim to relay. The host boundary prevents those invalid candidates from being treated as supported, but an unavailable packet is a failed workflow, not a correct engineering answer. This sequence makes source creation, candidate applicability, acquisition policy and recipient explanation separate parts of end-to-end correctness; the three amended synthetic schedules do not isolate a general EAL syntax benefit or establish a field effect.

Those false support events are outputs scored against a reference for the supplied information. The calls contain no contemporaneous human decision process and cannot be recoded as observed confirmation, automation or fluency bias in an engineer. Even if an agent flags a plausible bias mechanism, a wrong configuration, missing opportunity or strategic decision to economise on verification can yield the same visible outcome. Stage A supplies authored process records and explicit rivals, but its cases do not independently manipulate authority, order or presentation within a matched engineering record. It therefore tests reference agreement under those records rather than a human or model cue effect. The response boundary rejected 272 calls, including parsed answers with strong false attributions. Retaining those answers reveals failures hidden by the zero admitted false-attribution count. A future rationale review must also examine unsupported accusations in responses with weak structured status.

Several limits affect inference. The five diagnostic task templates and six relay templates were public development material, and related cases share mechanisms. The later packet pilots reused four roots and did not exercise current claim-scoped delivery. The new cross-model cohorts remain developmental: one reuses those exposed roots, while the other has three singly authored roots and lacks masked independent adjudication. Their equal checker has separate code, but selected endpoints were matched offline and supplied to both checked arms; 21/21 final statuses per model on either arm do not estimate recipient reasoning accuracy. Repetitions and evidence revisions are dependent; task-cluster intervals from six clusters have poor coverage in the repository's sensitivity simulation \citep{earlRelay}. A failed schema or output limit counts against an end-to-end answer but does not reveal the model's latent status judgement. Short public justifications and exploratory AI ratings are archived, yet no independently human-adjudicated explanation-quality estimate is available. A correct interpreter answer is conditional on its source, method contracts and reference fixtures, not a measurement of real plant behaviour. The new family resolver checks only finite tuples with a recorded review assertion, and its lexical retriever has ten single-author offline candidate queries and twelve singly authored model-selection tasks, with no independently adjudicated or field selection accuracy. The Stage A topology comparison additionally changes call allocation, and its model comparison changes reasoning settings. Its twelve authored families and five fully warranted scores cannot support a general model ranking or a calibrated population interval. The stopped identity diagnostic, request-route prompt conflict, outcome-informed Stage A continuation and post-hoc format readout further limit confirmatory interpretation. Repeated adaptive use of an exposed task set can reward local tuning and weaken a generalisation claim \citep{dwork2015}.

A controlled study of Toulmin structure on two source articles found a positive effect for one, with variable and difficult structure authoring \citep{cheikes2004}. This does not estimate EAL/2 effectiveness, but it makes a content-matched prose challenge necessary for the human comparison.

An independent test should sample new mechanisms across electronics and control, mechanical and thermal systems, distributed software, and reliability or manufacturing. Task authors should check the formal and engineering answers independently, including underdetermined cases. Match model class and tool access for EAL and the equal checker; score raw source, instruction-only, model-requested assessment, host-owned status and independent checker routes. Freeze task genealogy, model snapshots, prompts, budgets, status and explanation rubrics, source-review criteria and an information target before revealing outcomes. Measure candidate selection separately from exact family resolution. Retain failed calls, source revisions and human authoring/review effort. Such a study could support a model- and task-bounded causal delivery claim if assignment, measurement and coverage checks hold. A notation-specific claim additionally requires independently authored equal representations and comparable operations. External engineering cases would still be needed for field-use claims.

\section{Conclusion}\label{sec:conclusion}
EAL/2 supplies a finite, executable account of named argument dependencies, typed method results and targeted objections. Its claim-scoped host can keep an authorised checked status final despite inconsistent recipient prose, while finite family resolution and retrieval still require trusted applicability review. The historical raw-only comparison found EAL 4/10 against meaning-equivalent JSON 5/10; the new 96-call authoring schedule accepted no generated source and found one strictly faithful, exact one-shot explanation in 48 exploratory AI reviews. The 960-call information study recorded 130 admitted false supports, concentrated where invalid raw records were supplied. In the amended 800-slot deployment, fixed-source host and native routes delivered checked statuses by construction, but every final authored candidate failed validation and all 320 recipient packets were unavailable. Two exploratory AI reviews agreed on four faithful explanations among 48 selected responses, all from the raw route. A status result remains conditional on the authored warrant and acquisition chain: the Kubernetes contention example needs monitored coverage and mandatory reassessment as resources change. The finite tasks do not require an EAL grammar or ASPIC+ extension. The bias-attribution Stage A run retained 432 terminal outcomes but admitted only 98 of 288 route dispositions; five met the full reference. Twenty strong false attributions occurred in parsed answers rejected by subsequent validation. These records expose limits of the response contract and usable coverage; they establish neither safety from a zero admitted-error count nor a general topology or model advantage. The selected synthetic outcomes also leave language-specific benefit and live engineering truth unresolved. The separately frozen matched-format and equal-allocation plans have no provider results; the staged human-decision tool has no participants or outcomes. An independent human trial is required to test the original critical-analysis aim. Independent human adjudication, source-authoring repair tests and complete effort accounting are prerequisites for a meaningful all-in cost per correctly accepted and faithfully explained decision.

\section*{Data and code availability}
The historical relay and notation results use repository revision \texttt{2b0815f}; the later host, family, checker and offline schedules began at revision \texttt{d78836b} of \url{https://github.com/emmett08/earl}. The two 48-call pilot archives are under \path{benchmarks/results/2026-09-24-artifact-live-pilot/} and \path{benchmarks/results/2026-09-24-artifact-live-replication/}. The initial offline freeze and stopped attempt are under \path{benchmarks/results/2026-09-24-cross-model-offline/}; separate revised live archives, their read-only audits and a stricter recipient-status-consistency analysis are under \path{benchmarks/results/2026-09-24-cross-model-live/} (\path{stricter-readonly-analysis.tar.gz}). Their exact input SHA-256 values are \path{3d6a33d88720772f63255a1fac1210c05717b08fe4b49d7900a176920939529d} and \path{b8f9cf8fa03ac03c446c94d4e32287edbc0f45fcf4360c00195046375b8e13ef}. The corrected route's terminal ledgers and independent audits are under \path{benchmarks/results/2026-09-24-cross-model-route-v2/}; its complete 63-output exploratory review is under \path{benchmarks/results/2026-09-24-cross-model-route-v2-explanation-review/}. The separate de-identified 105-output cross-arm packet and AI reviews are under \path{benchmarks/results/2026-09-24-cross-model-explanation-review/}. The synthetic Kubernetes argument, application gate and offline family study are under \path{examples/kubernetes-resource-revision/} and \path{benchmarks/experiments/kubernetes-host-revisions/}. The distinct paid candidate-selection archive and its audit are under \path{benchmarks/results/2026-09-24-kubernetes-family-selection/}, frozen at SHA-256 \path{b5ef86e16f16d539d75d0632497b73862a6f6ed4c1ea9cd42d6fdb313244c7cc}. The archived notation analysis can be rerun without model calls using \path{scripts/analyse_notation_transfer.py} and \path{scripts/analyse_notation_format_sensitivity.py}. PR~7 merged the host, family, Kubernetes and earlier live records into \texttt{4c44e65b2aecbde6d9a0e60cc4290d4847134c0d}. The separately versioned negative-revision amendment, grammar decision and developmental campaign were merged in PR~8 at \texttt{b9177df}, with its 96-call freeze, raw ledger and explanation reviews under \path{benchmarks/results/2026-09-24-negative-revision-96-developmental/}. The original version 0.1.0 protocols are preserved. The amended 960-call archive is \path{benchmarks/results/eal2-960-campaign-20260924.tar.gz}, SHA-256 \path{6c55806f2e27cde1d4f2e7bc55d11893afc14e95309fed0578f7c10f2ba58900}; it retains the four immutable ledgers, composite and analysis. The fresh amended 800-slot archive is \path{benchmarks/results/2026-09-24-deployment-800-a3-developmental/a3-terminal.tar.gz}, SHA-256 \path{eabba61eae35a747fbb1fb209b06355ad12ad76126a4e3824b79b9706eab02a9}; its A1, A2 and C1 stopped diagnostic attempts remain in separate result directories and are not pooled with A3. The preselected A3 actual-prompt review frame and two independently frozen AI ratings are \path{a3-postcall-review-frame.tar.gz} and \path{a3-ai-review-records.tar.gz} in the A3 result directory; the latter has SHA-256 \path{0054a5394f04dad072e6aca39fd8a0404009ae6513cdda5bbee80fe0ad4fab31}.

The executed bias-attribution Stage A record was recovered at revision \path{430945cd24588240fc5ea351b723f4cfc924dea5} and integrated in PR~14 at \path{f2ba7ebbe3e37c6bb1e9c5cbbb02586f9ca06273}. The protocol, amendments 0.1.1 and 0.1.2, case manifest, collector and analyser are under \path{benchmarks/experiments/bias-mechanisms/}. The first stopped pilot is under \path{benchmarks/results/2026-09-24-bias-agent-pilot-stopped/}; completed pilot and full ledgers, reports, compressed freezes and historical commands are under \path{benchmarks/results/2026-09-24-bias-agent-stage-a/}. Pilot and full internal freeze digests are \path{71f160f0fcb28118180bb9b191e720cd95839dc40c44aec3cc4f8e523ea891a5} and \path{5c814040ee79627d844d05752ce5684b5984b8d67c2e6c60e87a133dec390c8e}. Their archive README records each file's SHA-256. The original zero-call plans remain historical; PR~13's alternative source-refresh proposal was closed without replacing executed material.

The paper's \path{verify_stage_a.py} reconstructs the pinned analysis material from local Git objects, verifies archive hashes and ledger chains, reproduces both reports byte for byte and generates \path{stage-a-summary.json} and \path{stage-a-table.tex}. It makes no provider calls. The source snapshot matters: later EAL modules and changed material hashes prevent reproduction by simply executing the historical freeze against current \texttt{main}. The script's \texttt{--check} option verifies the committed paper outputs.

The separate follow-up code, protocol, offline tests and execution note are under \path{benchmarks/experiments/bias-followups/}; its zero-call pilot and full schedules are under \path{benchmarks/results/2026-09-24-bias-followups-zero-call/}, at internal SHA-256 values \path{1d543e38827ee95ec7e17535ff6c1777d60e9a967d7677d1be696d88de05b678} and \path{1f790f06b2078a6244e9d7039b9a51f43e503df81dd71cc303e23d74dcdc68f2}. Their configured reserves are US\$2.0363 and US\$12.3126; actual charges are zero. The human allocation, packet and descriptive-scoring protocol and tests are under \path{benchmarks/experiments/bias-human-trial/}. No reviewed human case manifest, participant allocation, raw response or human outcome is archived. The newer independent-cohort plans are under \path{benchmarks/experiments/bias-next/} and were integrated through PR~11. These design records have no provider or participant outcomes. The repository is private; external reviewer access remains to be arranged by the author.

\section*{Funding}
The author must supply a factual funding statement before submission.

\section*{Declaration of competing interest}
The author must supply a factual competing-interest declaration before submission.

\section*{Author contributions}
The author must verify the work and supply a CRediT contribution statement before submission. An AI system is not an author.

\section*{Declaration of generative AI and AI-assisted technologies in the manuscript preparation process}
OpenAI ChatGPT (Codex) assisted with manuscript drafting, source checking, data-summary checking and LaTeX preparation. The human author must review and edit the content, verify every claim and reference, and take responsibility for the submitted article.

\begingroup
\raggedright
\bibliographystyle{elsarticle-num}
\bibliography{references}
\endgroup
\end{document}
